%%
%% trilogy_dag.tex  --  paper-grade logical DAG, XVIII--XX
%% Every edge A -> B: ``B requires A in its proof''.
%% Edges O4 -> H1 and O2 -> H2 now carry explicit mechanisms.
%%
\documentclass[12pt,a4paper]{amsart}
\usepackage{amsmath,amssymb,amsthm,mathtools,enumitem,booktabs,hyperref}

\newtheorem{theorem}{Theorem}[section]
\newtheorem{lemma}[theorem]{Lemma}
\newtheorem{proposition}[theorem]{Proposition}
\theoremstyle{definition}
\newtheorem{definition}[theorem]{Definition}
\newtheorem{remark}[theorem]{Remark}

\begin{document}

\title{Logical DAG: Trilogy XVIII--XX (Paper-Grade)}
\date{May 2026}
\maketitle

\tableofcontents

%% ================================================================
\section{Purpose and Conventions}
\label{sec:purpose}
%% ================================================================

This document is the \emph{paper-grade} logical dependency graph
for the trilogy XVIII--XX.
It is designed to withstand reviewer scrutiny: every edge is
explicit and every open edge carries a constructive mechanism
(not just a label).

\medskip
\textbf{Conventions.}
(T) = proved theorem in the trilogy.
(H) = hypothesis (axiom within $\mathfrak{T}_\gamma$).
(S) = proof sketch (power family).
(O) = open problem.
An edge $A\to B$ means B's proof invokes A.
An edge $A\xrightarrow{\text{mech}}B$ additionally names
the explicit mechanism.

%% ================================================================
\section{Layer 0: $\gamma$-Agnostic Foundations}
\label{sec:L0}
%% ================================================================

\begin{center}
\renewcommand{\arraystretch}{1.3}
\begin{tabular}{l l l l}
\toprule
\textbf{ID} & \textbf{Statement} & \textbf{Source} & \textbf{Used in} \\
\midrule
F1 & $S_K - K_{\mathcal{A}} = A_K - B_K$ & XVIII, HS theory
  & A5, U6 \\
F2 & $A_K = A_K^{\rm lin} + A_K^{\rm quad}$ & XIX, F1
  & U2 \\
F3 & $\|f\otimes g\|\leq M\|f\|_{C^0}\|g\|_{C^0}$ & XVIII, standard
  & U2 \\
F4 & Bessel tail & XVIII & A5 \\
F5 & Tonelli/DCT & XVIII, standard & A5 \\
F6 & Slepian tail $\|B_K\|\leq C_2^{1/2}e^{-\alpha K/2}$ & XVIII
  & A5, U5, U6 \\
F7 & Sequential $(c,K)$-limit & XVIII, F5, F6 & A5 \\
F8 & Kato norm-resolvent & XVIII, Kato VI & A5 \\
F9 & Total order $\mathcal{S} = \{c^{-s}(\log c)^t\}$
  & XX Def.~2.1, Hardy--Littlewood
  & U3, U4, U7 \\
\bottomrule
\end{tabular}
\end{center}

%% ================================================================
\section{Layer 1: Airy-Specific Results (XVIII--XIX)}
\label{sec:L1}
%% ================================================================

\begin{center}
\renewcommand{\arraystretch}{1.3}
\begin{tabular}{l l l l}
\toprule
\textbf{ID} & \textbf{Statement} & \textbf{Source} & \textbf{Used in} \\
\midrule
A1 & $a_j\sim(3\pi j/2)^{2/3}$
  & XVIII, Airy zero asymptotics
  & A3 \\
A2 & Langer approx $\psi_j\approx\mathrm{Ai}(c^{2/3}\zeta_j)$
  & XVIII, Langer (1949)/Olver (1954)
  & A3, A4 \\
A3 & $C_j = O(j^{2/3})$: Olver remainder
  & XVIII \texttt{rem:langerkdep}, A1, A2
  & A4, A6 (instantiates H1) \\
A4 & Turning-point window $c^{-1/3}$
  & XVIII, A2, A3
  & A5 \\
A5 & BR3: $\|K_{B_c}-K_{\mathcal{A}}\|\to 0$
  & XVIII (T), F1--F8, A2--A4
  & (qualitative baseline) \\
A6 & $C_{\rm lin}=O(K^{5/3})$, $C_{\rm quad}=O(K^{7/3})$
  & XIX (T), F2, F3, A3, summation
  & A7, A8 \\
A7 & $K^*_{\rm Airy}\sim 2/(5\alpha)\log c$
  & XIX (T), A6, F6,
    Lem.~\ref{lem:kopt_variational} (Airy)
  & A8 \\
A8 & Airy rate $O(c^{-1/3}(\log c)^{5/3})$
  & XIX (T), A6, A7, F6, F9
  & (instantiates U7 for $\gamma=2/3$) \\
\bottomrule
\end{tabular}
\end{center}

%% ================================================================
\section{Layer 2: Universal Results (XX, Conditional)}
\label{sec:L2}
%% ================================================================

\begin{center}
\renewcommand{\arraystretch}{1.3}
\begin{tabular}{l p{5.0cm} l l}
\toprule
\textbf{ID} & \textbf{Statement} & \textbf{Source} & \textbf{Req.} \\
\midrule
H1 & $\|e_j\|_{C^0}\leq C_0 j^\gamma c^{-\beta}$
  (pointwise Langer)
  & (H); A3 verifies $\gamma=2/3$;
    O4 closes for power family
  & O4 (S) \\
H2 & $C_\infty = \sup_j\|u_j\|_{C^0}<\infty$
  & (H); standard for Airy;
    WKB normalization for power family
  & O2 (S) \\
U1 & $\sum_{j=1}^K j^{n\gamma}\sim K^{1+n\gamma}/(1+n\gamma)$
  & XX (T), Euler--Maclaurin
  & --- \\
U2 & $C_{\rm lin}=O(K^{1+\gamma})$,
    $C_{\rm quad}=O(K^{1+2\gamma})$
  & XX Lem.~5.1 (T), F2, F3, H1, H2, U1
  & H1, H2 \\
U3 & Layer dominance:
    $c^{-n\beta}(\log)^{1+n\gamma}\prec c^{-\beta}(\log)^{1+\gamma}$
  & XX Prop.~2.2 (T), F9
  & F9 \\
U4 & Layer-Survival Theorem (all $n\geq 2$ subleading)
  & XX Thm.~7.1 (T), U2, U3
  & U2, U3 \\
U5 & $K^* = \mathrm{argmin}\,F(K)$;
    $K^*\sim 2\beta/(\alpha(1+\gamma))\log c$
  & XX Lem.~6.1 (T), F6, U2
  & F6, U2 \\
U6 & Two-scale bound
    $\|K_{B_c}-K_\mathcal{A}\|\leq C_{\rm lin}c^{-\beta}
    +C_{\rm quad}c^{-2\beta}+C_2^{1/2}e^{-\alpha K/2}$
  & XX Thm.~8.1 (T), F1, F6, U2
  & F1, F6, U2 \\
U7 & $O(c^{-\beta}(\log c)^{1+\gamma})$ at $K=K^*$
  & XX Cor.~8.2 (T), U4, U5, U6
  & U4, U5, U6 \\
\bottomrule
\end{tabular}
\end{center}

%% ================================================================
\section{Layer 3: Open Problems with Explicit Mechanisms}
\label{sec:L3}
%% ================================================================

\begin{center}
\renewcommand{\arraystretch}{1.3}
\begin{tabular}{l p{4.0cm} l p{4.5cm}}
\toprule
\textbf{ID} & \textbf{Statement} & \textbf{Closes} & \textbf{Mechanism} \\
\midrule
O4 & Complete Olver-Watson
  expansion for $Q_j\leq D_1 j^\gamma$
  in \texttt{olver\_growth\_lemma.tex}
  & H1 (power family)
  & Olver~\cite{Olver1997} Ch.~12 remainder
    $+$ turning-point WKB scaling:
    $Q_j = \int|\Psi(\zeta_j)|\,d\zeta$;
    $\xi_j\sim j^{\gamma/p}$ gives
    $Q_j\sim D_p j^{\gamma/p}\leq D_1 j^\gamma$
    (Airy worst case, $p=1$) \\
O2 & Verify $C_\infty<\infty$
  for $V=|\xi|^\nu$
  & H2 (power family)
  & WKB $L^\infty$ normalization:
    $\|u_j\|_{C^0}\leq C\cdot a_j^{-1/4}|V'(\xi_j)|^{-1/6}$;
    check growth in $j$ \\
O1 & Prove H1 for general $V$
  (beyond power family)
  & H1 (general)
  & Regular variation condition
    $\xi V'(\xi)/V(\xi)\to p$
    extends O4 mechanism to
    all r.v.\ potentials \\
O3 & Lower bound $\Omega(K^{1+\gamma})$
  for $C_{\rm lin}$ (sharpness)
  & Optimality of U4
  & Spectral counting + turning-point
    density argument \\
\bottomrule
\end{tabular}
\end{center}

%% ================================================================
\section{Critical Path Analysis}
\label{sec:critical_path}
%% ================================================================

\textbf{Minimal path to unconditional theorem (power family):}
\[
  \text{O4}\xrightarrow{\text{Olver Ch.12}+Q_j}
  \text{H1}\xrightarrow{}
  \text{U2}\xrightarrow{}
  \text{U4}\xrightarrow{}
  \text{U7}.
\]

\textbf{Independent path:}
$\text{O2}\xrightarrow{\text{WKB }L^\infty}\text{H2}\to\text{U2}$.

\textbf{Full generalization:}
$\text{O1}\xrightarrow{\text{reg.var.}}\text{H1 (general)}\to$
all of U2--U7 unconditionally for regularly varying $V$.

\medskip
\textbf{Critical observation.}
The O4$\to$H1 edge is no longer a bare label.
The mechanism is:
\begin{enumerate}[label=(\roman*),nosep]
  \item Olver Ch.~12 gives $\|e_j\|\leq 2AQ_j c^{-\beta}$
    with explicit remainder integral $Q_j$.
  \item WKB turning-point scaling gives
    $Q_j\sim D_p j^{\gamma/p}\leq D_1 j^\gamma$
    (Airy worst case).
  \item Therefore $C_0 = 2AD_1$ is independent of $j$, $K$, $c$.
\end{enumerate}
This is the constructive mechanism that was previously implicit.

%% ================================================================
\section{Dependency Graph (Textual)}
\label{sec:textual_dag}
%% ================================================================

\begin{verbatim}
[Classical: Olver Ch.12, Langer 1949]
     |                   |
     v                   v
   [A2] Langer approx  [A1] Airy zeros
     |                   |
     +--------+----------+
              |
              v
            [A3] C_j=O(j^{2/3})  ---instantiates---> [H1 (Airy)]
              |                                           |
              v                             [O4: Olver-Watson]
            [A5] BR3                         |  Q_j <= D_1 j^gamma
                                             |
                          [H1 (power)]<------+   [H2: C_inf]
                               |                    |
     [F2,F3,U1: summation]     |                    |
              |                +--------+-----------+
              |                         |
              v                         v
     [F9: scale S]              [U2: C_lin, C_quad]
              |                         |
       [U3: dominance]           [F6: Slepian tail]
              |                         |
              +------------+------------+
                           |
                     [U4: Layer-Survival]
                     [U5: K* variational]
                     [U6: two-scale bound]
                           |
                           v
                     [U7: O(c^{-beta}(log c)^{1+gamma})]  QED
\end{verbatim}

%% ================================================================
\section{Freeze Checklist}
\label{sec:freeze}
%% ================================================================

\begin{center}
\renewcommand{\arraystretch}{1.3}
\begin{tabular}{l l}
\toprule
\textbf{Item} & \textbf{Status} \\
\midrule
All F1--F9 in XVIII or XX & \checkmark \\
All A1--A8 proved in XVIII--XIX & \checkmark \\
H1, H2 explicitly stated as hypotheses & \checkmark \\
U1--U7 proved conditional on H1, H2 & \checkmark \\
O1--O4 carry explicit mechanisms & \checkmark \\
$K^*$ via variational Lemma~6.1; FOC~\eqref{eq:FOC}
  & \checkmark \\
$c$-scale comparisons via Def.~2.1 (total order) & \checkmark \\
O4$\to$H1 edge constructive (Olver$+Q_j$ mechanism)
  & \checkmark \\
Conditional theorem scope in abstract & \checkmark \\
H1 status table per model family & \checkmark \\
\hline
O4 (Olver-Watson complete proof) & open \\
O2 ($C_\infty$ for $\nu\neq 1$) & open \\
\bottomrule
\end{tabular}
\end{center}

\medskip
\textbf{Freeze verdict.}
XX is paper-grade ready for freeze as a conditional theorem.
The two open items (O2, O4) are self-contained sub-problems;
they do not affect internal consistency or the correctness
of any proof given H1--H2.
A future Paper~XXI could address O1 (regular variation)
and O3 (lower bound) as part of a rigorous closure programme.

\begin{thebibliography}{99}
\bibitem{Olver1997} F.~Olver,
  \textit{Asymptotics and Special Functions}, AK Peters, 1997.
\bibitem{Langer1949} R.~Langer,
  Trans.\ Amer.\ Math.\ Soc.\ \textbf{67} (1949), 461--490.
\bibitem{Olver1954} F.~Olver,
  Phil.\ Trans.\ Roy.\ Soc.\ London \textbf{247} (1954), 307--327.
\end{thebibliography}

\end{document}
